\clearpage % TEMP
\section{Preliminaries}\label{pre}
% ------------------------------------------------------------------------------
\paragraph{Assumptions}
For simplicity, we first assume that
periods of sex work have clear start/end points,
and that formal sex work can be clearly distinguished from transactional sex.
We also assume that surveys predominantly reach those actively engaged in sex work,
despite some survey eligibility criteria including
individuals who sold sex anytime in the past year, etc. \cite{Scorgie2012}.
We further assume that survey participants
can accurately report how long they have sold sex
--- \ie perfect recall, no social desirability bias.
Finally, we assume that all populations are at equilibrium,
\eg rates of entering/exiting sex work are stable over time.
% ------------------------------------------------------------------------------
\paragraph{Who exactly are we talking about?}
Before continuing, it is helpful to distinguish two populations:
individuals who will ever sell sex in their lifetime, and
individuals who are actively selling sex.
The former "source population"
is the target sample for method~1 in \cite{Fazito2012},
and tends to include more individuals who sell sex for shorter durations;
the latter "active population"
(a subset of the source population)
is the target sample for methods 2 and 3 in \cite{Fazito2012},
and tends to include more individuals who sell sex for longer durations.
Figure~\ref{fig:diag} illustrates periods in sex work
for a hypothetical source population at equilibrium, distinguishing
the active population (green) from the inactive population (gray)
at the time of a hypothetical survey (yellow).
% LW: I wonder if you can elaborate more such that
%     Figure 1 can be supportive of the statement more explicitly
%     (it is doing the work but i think readers may not see it directly):
%     e.g., u can highlight that the orange line (survey) - which is often the latter sampling
%     often capture more dark green - which tend to have longer duration selling/buying)
%     and omit the grey (those who had finished selling sex and tend to have shorter duration) ...
% JK: I added a more detailed reference to the figure in a new sentence,
%     rather than trying to integrate references to the Figure
%     into the existing sentences, since they are already somewhat information-dense,
%     (the new label, longer vs shorter durations, and referencing methods 1,2,3).
% JK: I'm not sure if we need to explain the equilibrium thing more,
%     i.e. that each period is a new person, not multiple periods by the same person,
%     but that the equilibrium assumption implies people will tend to get "replaced"
%     by someone with similar duration during turnover.
\par
Since "sex work" strata in compartmental models
represent individuals actively selling sex,
we might assume that rates of exit from these strata should be derived from
durations among the active population.
However, such rates should actually be derived from
durations among the source population,
because exit rates apply to all individuals \emph{entering} sex work,
which reflects the source population.
% LW: I think this may need further clarification. Sth along
%     "Compartmental models of HIV often include a 'sex work' strata
%     which captures the period when an individual/risk group actively sells sex.
%     Entry into and exit from this strata are modeled using estimates on
%     rates of entering sex work and rates of existing from sex work.
%     Although the strata captures active sex work, the rates of enter and exit
%     however apply to the source population (who will ever sell sex).
In an analogous example:
the duration of an infection should be characterized by
a cohort of individuals with \emph{incident} infections,
not a cross-section of individuals with \emph{prevalent} infections.
\par
A final population worth distinguishing is
the "sampled population" in a given cross-sectional survey,
which we assume intends to be representative of the active population,
but might undersample some individuals,
\eg those in the early years of sex work,
due to delays in self-identification as a sex worker
\cite{Mishra2012,Cheuk2020}.
\begin{figure}
  \centering\includegraphics[scale=1]{diag.leg}
  \caption{Diagram of periods selling sex in an open population at equilibrium}
  \label{fig:diag}
  \floatfoot{Note: consecutive periods in each row denote multiple individuals,
    not multiple periods of sex work by the same individual.}
\end{figure}
